% !TEX program = xelatex
\documentclass[11pt,a4paper]{ctexart}
\usepackage{amsmath,amssymb,amsfonts}
\usepackage{enumitem}
\usepackage{booktabs}
\usepackage{geometry}
\geometry{margin=2.5cm}

\title{\texttt{chen\_greedy\_algorithm}：每一步怎么定「搬哪个箱、搬到哪」}
\author{上港堆场整理 --- \texttt{greedy} 包}
\date{}

\newcommand{\capG}{G_{\mathrm{上限}}}

\begin{document}
\maketitle

\section{一句话}

算法\textbf{反复做同一件事}：在当前堆场上，挑\textbf{一个}合法搬箱动作——从某堆\textbf{顶}上取走一箱，放到\textbf{另一堆顶}上——执行后进入下一步，直到要么全场排好，要么再也挑不出合法动作。

\textbf{关键：}每一步\textbf{只动顶箱}，且\textbf{不能同堆内倒层}（只能从堆 $i$ 搬到堆 $k$，$i\neq k$）。

\section{每一步的六步流程（白话）}

下面用「当前已经搬了 $m$ 步」时的状态说明。记全场逆序箱数为 $N$，初始为 $N^{(0)}$，允许的上限为
\[
  \capG = \bigl\lfloor 1.3 \times N^{(0)} \bigr\rfloor,
  \qquad
  G = m + N \quad \text{（搬箱步数 + 当前逆序箱数）}.
\]

\begin{enumerate}[label=\textbf{第 \arabic* 步}, leftmargin=2.2em]

  \item \textbf{先看哪些堆「值得动」}

  只从下面两类堆里往外搬箱（源堆）：
  \begin{itemize}
    \item 这堆还有逆序（$N_i>0$），或
    \item 这堆还没排好（自下而上关序号不是「下面大、上面小」）。
  \end{itemize}
  其它已经排好的空问题堆，这一步\textbf{不会}从它往外搬。

  \item \textbf{定「搬哪个箱」——只能是源堆最顶上那一箱}

  对每个源堆 $i$，没有选择权：\textbf{被搬的箱已经定了}，就是堆 $i$ 最上面那一箱（顶箱）。
  程序里记它的秩为 $n$，源层号 $j=|S_i|$（当前堆高）。

  \item \textbf{定「能搬到哪」——枚举所有别的堆的顶}

  对源堆 $i$，逐个看目标堆 $k$（$k\neq i$）：
  \begin{itemize}
    \item 若堆 $k$ 已满（层数 $=H$），跳过；
    \item 否则允许：把该顶箱放到堆 $k$ 的\textbf{新顶层}（层号 $\ell=|S_k|+1$）。
  \end{itemize}
  于是得到一个候选：\textbf{「堆 $i$ 的顶箱 $\to$ 堆 $k$ 的顶」}。
  所有源堆 $\times$ 所有能接箱的目标堆，组成这一步的\textbf{候选动作表}。

  \item \textbf{对每个候选先「试搬」再打分，不合格的划掉}

  对表里每一条「$i$ 顶箱 $\to$ $k$ 顶」，在内存里复制当前布局、模拟搬一次（不真改现场），然后检查：
  \begin{enumerate}[label=(\alph*)]
    \item 搬完之后 $G' = (m+1) + N'$ 是否 $\le \capG$？超过则\textbf{淘汰}（翻箱代价超标）；
    \item 搬完后的整盘布局是否\textbf{以前出现过}？若重复且已有更低累计逆序代价，则\textbf{淘汰}（防止来回打转）。
  \end{enumerate}
  剩下的就是\textbf{合法候选}。

  \item \textbf{在合法候选里挑「最好」的一条——这就是本步真正执行的动作}

  对每条合法候选算以下指标：
  \begin{itemize}
    \item $\sum_{\tau=0}^{t}N_\tau$：截至当前路径的累计未来逆序箱数（主目标，优先\textbf{尽量小}）；
    \item $N'$：搬完后的当前未来逆序箱数（次目标，优先\textbf{尽量小}）；
    \item $m+1$：搬箱步数（再次级目标，优先\textbf{尽量小}）；
    \item 多贝时加分 $B$：源堆在初始逆序最多的贝（\texttt{drain\_bay}）$+2$，目标堆在初始逆序最少的贝（\texttt{fill\_bay}）$+1$（在仍打平时优先\textbf{大}）。
  \end{itemize}

  \textbf{排序规则（从上到下比，先比赢的优先）：}
  \[
    \sum_\tau N_\tau \ \text{小}
    \ \rightarrow\
    N' \ \text{小}
    \ \rightarrow\
    m+1 \ \text{小}
    \ \rightarrow\
    B \ \text{大}.
  \]
  取排在最前面的一档；若这一档还有多条，\textbf{随机抽一条}执行。

  \item \textbf{真搬、记步、进入下一轮}

  把选中的「堆 $i^\ast$ 顶箱 $\to$ 堆 $k^\ast$ 顶」写入动作序列，更新现场堆场，$m\leftarrow m+1$，把新布局记入「已出现」集合。
  若此时每一堆都已良序，\textbf{结束}；否则回到第 1 步。

\end{enumerate}

\section{小例子（帮助对齐直觉）}

假设当前只有两堆有问题：堆 $3$ 顶箱关序号偏大（有逆序），堆 $7$ 还能再接一箱。
\begin{itemize}
  \item \textbf{搬哪个箱？} 若 $3\in\mathcal{I}_{\mathrm{src}}$，本步只可能搬\textbf{堆 3 顶箱}，不能搬堆 3 下面的箱。
  \item \textbf{搬到哪？} 若堆 $7$ 未满，则产生候选「$3$ 顶 $\to$ $7$ 顶」；若堆 $5,7,9$ 都能接，则产生三条候选，分别试搬、过滤、再比 $\Delta N$ 和 $G'$。
  \item \textbf{为何可能不搬 $3\to 5$？} 试搬后若 $N$ 反而变大，或 $G'$ 超过 $\capG$，或布局重复，该候选会被划掉，最终可能只剩 $3\to 7$ 合法。
\end{itemize}

\section{指标（需要时再查公式）}

\subsection{什么叫「排好」}

从一堆\textbf{底}往\textbf{顶}看关序号：每一层都要比上面一层\textbf{更大}（先提的箱关序号更小）。满足则该堆良序。

\subsection{逆序箱数 $N_i$（单堆）}

从第 2 层往上找\textbf{第一处}「上面关序号 $>$ 下面关序号」的层 $j^\ast$；从 $j^\ast$ 到堆顶所有非空箱的个数，即该堆 $N_i$。全场 $N=\sum_i N_i$。

\subsection{为何用 $G=m+N$}

每搬一步 $m$ 加 1，逆序 $N$ 可能变；要求任意时刻 $m+N \le \lfloor 1.3 N^{(0)}\rfloor$，控制「搬箱次数 + 剩余乱度」的总代价。搜索排序时以累计未来逆序箱数$\sum_t N_t$为主目标，而不是硬性要求每一步$N$都下降。

\section{何时停}

\begin{itemize}
  \item \textbf{成功：} 所有堆都良序；
  \item \textbf{卡住：} 第 4 步后合法候选为空（在满足 $G\le\capG$ 的前提下无路可走）；
  \item \textbf{保底：} 步数达到程序上限 \texttt{hard\_cap}。
\end{itemize}

\section{和代码的对应}

\begin{table}[h]
  \centering
  \begin{tabular}{p{5.5cm}p{6.5cm}}
    \toprule
    白话步骤 & 函数 \\
    \midrule
    哪些堆当源堆 & \texttt{greedy\_plan.\_src\_stacks} \\
    顶箱搬到哪些堆顶 & \texttt{greedy\_plan.\_enumerate\_cross} \\
    试搬 + 三条过滤 & \texttt{simulate\_move} + \texttt{\_pick\_best\_move} \\
    比累计 $N$、当前 $N$、步数和 $B$ 并随机 tie & \texttt{\_pick\_best\_move} 内排序 \\
    真搬并循环 & \texttt{apply\_move} + \texttt{plan\_greedy\_until\_terminal} \\
    \bottomrule
  \end{tabular}
\end{table}

\end{document}
